\section{Related Work}
\label{sec:related-work}

\paragraph{Event-level evaluation.}
Event detection and changepoint detection have long been treated separately from samplewise classification because a long sequence may matter primarily through the events it contains \citep{guralnik1999,aminikhanghahi2017}.
Modern benchmarks increasingly use the ranked-matching structure of object and activity detection: predicted times are sorted by confidence, matched to annotations under a tolerance rule, and summarized by precision--recall or average precision \citep{everingham2010pascal,lin2014coco}.
SoftED and range-aware precision--recall make the same point for time series, where per-sample accuracy can disagree with event-level utility \citep{salles2023,tatbul2018precision}.
This paper follows that evaluation view back to the training objective.
If the metric ranks matched events, the training signal should describe the events that can be matched in a time region, rather than only which state label holds at each sample.

\paragraph{Segmentation and localized supervision.}
Samplewise segmentation models remain practical because they share computation across long recordings, and segmentation architectures are well established in sleep and biomedical time series \citep{supratak2017,perslev2019,phan2019,ronneberger2015}.
The relevant choice is the loss attached to such sequence predictors when the benchmark consumes detections.
Pose estimation, acoustic-event detection, and temporal action localization provide useful precedents for placing supervision near annotated times through keypoint maps, proximity targets, or start/end evidence \citep{newell2016stacked,xiao2018simple,luo2021,yu2022,phan2015,gao2017turn,lin2018bsn,lin2019bmn}.
Recent time-series event work makes a similar criticism of samplewise event/non-event labels, especially for rare and interval-defined events, and argues for a broad regression-based alternative to classification \citep{azib2023universal}.
These precedents establish localized event targets as prior work.
They motivate moving supervision toward the times that will be detected, but a proximity target alone does not specify what the output should estimate after the recording is windowed, smoothed, or downsampled.
BDL addresses the complementary question of what that localized supervision estimates once a continuous recording is clipped, smoothed, and mapped to a shorter output timeline.
If localized labels are arbitrary proximity scores, their scale can depend on window length, kernel width, pooling, or nearby annotations.
BDL fixes the unit and the scoring rule instead: each annotation contributes one event, smoothing distributes that contribution over plausible times, binning sums contributions onto the model timeline, and a proper log score fits the conditional mean.
The event unit therefore keeps the same interpretation before and after discretization, including windows with no events or multiple nearby events.

\paragraph{Counting-process likelihoods.}
Temporal point processes provide the standard language of counting measures, conditional intensities, and likelihoods over event times \citep{kingman1993,du2016rmtpp,mei2017neural,shchur2021neural}.
BDL borrows the counting-measure view for supervised prediction: given an observed window, it predicts the expected number of scored events in each output bin without modeling censoring, future histories, survival terms, or mark distributions.
The Poisson negative log-likelihood is a proper log score for this conditional mean \citep{gneiting2007proper}, making it a natural objective when the benchmark consumes temporal detections rather than complete event histories.
This places BDL between segmentation and full event-history modeling.
It keeps the efficient sequence-prediction interface used by segmentation models, while giving each output the likelihood interpretation of expected events per bin.
